\subsection{用位刻画整元素}

命题 5. 设 K 是域，S 是 K 的子环，h 是从 S 到某个域的同态，p 是 h 的核。为使 K 的元素 x 在局部环 \(S_{p}\) 上整，必须且只需 K 的每个延拓 h 的位都在 x 处有限。

若 \(f\) 是 \(K\) 的延拓 \(h, f\) 在 \(S_p\) 上有限、在 \(pS_p\) 上为零，从而位 \(f\) 的环支配 \(S_p\)。反之，若 \(V\) 是 \(K\) 的支配 \(S_p\) 的赋值环，则 \(V\) 是某个位 \(f\) 的环，它在 \(S\) 上的限制是与 \(h\) 有相同核的同态；把 \(f\) 换成一个等价的位，即可看出 \(V\) 是 \(K\) 的某个延拓 \(h\) 的位的环。说 \(K\) 的每个延拓 \(h\) 的位都在 \(x\) 处有限，等价于说 \(x\) 属于 \(K\) 的支配 \(S_p\) 的所有赋值环。于是命题由 § 1 第 3 节定理 3 得出。

命题 6. 设 K 是域，S 是 K 的子环。为使元素 \(x \in \mathbf{K}\) 在 S 上整，必须且只需 K 的每个在 S 上有限的位都在 \(x\) 处有限。

这同样是 §1 第 3 节定理 3 的推论。

